\newpage
\begin{center}
{\fontsize{16pt}{19.2pt}\selectfont \textbf{XÂY DỰNG MÔ HÌNH HỌC TĂNG CƯỜNG TỐI ƯU HÓA CHÍNH SÁCH ĐẶT HÀNG LẠI TRONG QUẢN LÝ TỒN KHO ĐA KHO}\par}
\vspace{0.5cm}
{\fontsize{16pt}{19.2pt}\selectfont \textbf{TÓM TẮT}\par}
\end{center}
\vspace{0.2cm}

Quản lý tồn kho trong mạng lưới bán lẻ nhiều kho, nhiều mặt hàng là bài toán ra
quyết định tuần tự dưới bất định: phải cân bằng chi phí lưu kho, thiếu hàng, đặt
hàng và tràn kho trong khi vẫn giữ mức phục vụ khách hàng. Các chính sách cổ điển
như EOQ, $(s, S)$ và Newsvendor dựa trên giả định nhu cầu dừng, thời gian giao
hàng cố định và không ràng buộc sức chứa, nên khó giữ tính tối ưu khi các giả
định này bị vi phạm.

Khóa luận mô hình hóa $10$ kho $\times$ $30$ mặt hàng thành $300$ tác tử, mỗi
tác tử quyết định lượng đặt hàng cho một cặp kho--mặt hàng từ quan sát cục bộ,
và các tác tử trong cùng kho chia sẻ một sức chứa hữu hạn. Các tác tử được huấn
luyện bằng IPPO với chia sẻ tham số (PS-IPPO), dùng phần thưởng phân rã cục bộ
theo từng cặp, chuẩn hóa theo quy mô chi phí của từng cặp, và mục tiêu mức phục
vụ $85\%$ triển khai dưới dạng ràng buộc mềm.

Môi trường mô phỏng dùng nhu cầu thực tế của bộ dữ liệu M5 Walmart ($1.941$
ngày), chia thành ba miền huấn luyện, validation và kiểm tra, với thời gian giao
hàng ngẫu nhiên từ một đến ba ngày. PS-IPPO được so sánh với ba chính sách cổ
điển đã tinh chỉnh trên miền validation, qua nhiều hạt giống huấn luyện và kiểm
định thống kê theo cặp.

% CAP NHAT SAU DOT CHAY CUOI: cac con so duoi day tu ban 3 hat giong x 1.000 episode.
Kết quả cho thấy khi được so ở cùng mức phục vụ, PS-IPPO đạt tổng chi phí vận
hành thấp hơn $(s, S)$ $10{,}2\%$ và thấp hơn Newsvendor $14{,}4\%$, còn EOQ không
đạt được mức phục vụ đó. Ưu thế tập trung ở các giai đoạn nhu cầu biến động
(ngày có sự kiện, cuối tuần); khi nhu cầu ổn định, $(s,S)$ đạt chi phí thấp hơn.
Chia sẻ tham số duy trì hiệu năng tốt ở nhóm mặt hàng có nhu cầu trung bình và
cao, nhưng phân bổ mức phục vụ thấp hơn cho nhóm nhu cầu thấp. Đóng góp của khóa
luận là một khung PS-IPPO có ràng buộc sức chứa dùng chung trên dữ liệu bán lẻ
thực tế, cùng một giao thức đánh giá tách biệt rõ chi phí vận hành, mức phục vụ
và tín hiệu huấn luyện.

\vspace{0.3cm}
\noindent\textbf{Từ khóa:} quản lý tồn kho; học tăng cường đa tác tử; IPPO; chia
sẻ tham số; M5.
